\documentclass[11pt]{article}
\usepackage[margin=0.8in]{geometry}
\usepackage{amsmath,amssymb,amsthm}
\usepackage{enumitem}
\usepackage{longtable}
\usepackage{array}
\usepackage{xcolor}
\usepackage{hyperref}
\usepackage{fancyhdr}
\usepackage{titlesec}
\usepackage{tcolorbox}
\usepackage{multicol}
\usepackage{booktabs}
\hypersetup{colorlinks=true, linkcolor=blue!50!black, urlcolor=blue!50!black}
\pagestyle{fancy}
\fancyhf{}
\rhead{\thepage}
\lhead{Engineering Math to Pure Math}
\setlength{\parskip}{0.55em}
\setlength{\parindent}{0pt}
\setlist[itemize]{topsep=0.2em,itemsep=0.2em}
\setlist[enumerate]{topsep=0.2em,itemsep=0.25em}
\newtcolorbox{focusbox}{colback=blue!3!white,colframe=blue!50!black,title=Focus}
\newtcolorbox{intuitionbox}{colback=green!3!white,colframe=green!40!black,title=Intuition}
\newtcolorbox{formalbox}{colback=gray!4!white,colframe=gray!70!black,title=Formal}
\newtcolorbox{walkbox}{colback=orange!5!white,colframe=orange!60!black,title=Walking question}
\newcommand{\N}{\mathbb{N}}
\newcommand{\Z}{\mathbb{Z}}
\newcommand{\Q}{\mathbb{Q}}
\newcommand{\R}{\mathbb{R}}
\newcommand{\C}{\mathbb{C}}
\newcommand{\eps}{\varepsilon}
\newcommand{\pow}{\mathcal{P}}
\theoremstyle{definition}
\newtheorem{definition}{Definition}
\newtheorem{example}{Example}
\newtheorem{exercise}{Exercise}
\theoremstyle{plain}
\newtheorem{theorem}{Theorem}
\newtheorem{lemma}{Lemma}
\theoremstyle{remark}
\newtheorem{remark}{Remark}

\title{\textbf{Proof Debugger}\\A Workflow for Finding and Repairing Gaps}
\author{Engineering Math to Pure Math}
\date{May 2026}
\begin{document}
\maketitle

\section{The Debugger Mindset}

A proof is not a poem where the mood carries the meaning. It is closer to a program with types, inputs, outputs, and invariants. But unlike a program, the compiler is another human mind. Your job is to make every transition checkable.

\section{Proof Submission Template}

\begin{enumerate}
\item \textbf{The claim:} State the theorem exactly.
\item \textbf{Quantifiers:} Write all hidden quantifiers.
\item \textbf{Definitions:} List the definitions you expect to use.
\item \textbf{Strategy:} Direct, contrapositive, contradiction, cases, induction, construction, or equality by two inclusions.
\item \textbf{Proof attempt:} Write the proof.
\item \textbf{Self-debug notes:} Mark the lines you are least sure about.
\end{enumerate}

\section{Feedback Rubric}

\begin{longtable}{p{0.2\linewidth}p{0.36\linewidth}p{0.34\linewidth}}
\toprule
Issue type & Symptom & Repair \\
\midrule
Quantifier gap & You prove something for one object but need all objects & Start with ``Let $x$ be arbitrary'' or produce the needed witness \\
Definition gap & You use a term without unpacking it & Replace the term by its definition \\
Implication gap & A line does not follow from the previous line & Add a lemma or algebraic justification \\
Case gap & You handle one possibility but not all & List exhaustive cases before proving them \\
Existence gap & You say something exists without constructing or citing it & Name the object and verify its properties \\
Uniqueness gap & You show there is one object but not only one & Assume two candidates and prove they are equal \\
Notation gap & Symbols change meaning mid-proof & Declare domains and variables before using them \\
\bottomrule
\end{longtable}

\section{Example Debug}

\textbf{Claim.} If $n^2$ is even, then $n$ is even.

\textbf{Weak attempt.} Since $n^2$ is even, $n^2=2k$. Taking square roots, $n=\sqrt{2k}$, so $n$ is even.

\textbf{Diagnosis.} The last step is not justified. Being the square root of twice an integer does not mean being twice an integer.

\textbf{Repair by contrapositive.} We prove: if $n$ is odd, then $n^2$ is odd. Suppose $n$ is odd. Then $n=2k+1$ for some $k\in\Z$. Hence
\[
n^2=(2k+1)^2=4k^2+4k+1=2(2k^2+2k)+1.
\]
Since $2k^2+2k\in\Z$, $n^2$ is odd. Therefore, by contrapositive, if $n^2$ is even, then $n$ is even.

\section{Line-by-Line Checklist}

\begin{enumerate}
\item Did every variable get introduced?
\item Did every introduced variable have a domain?
\item Did the proof use the exact hypothesis?
\item Did the conclusion match the theorem exactly?
\item Did any sentence hide the real work behind ``clearly,'' ``obviously,'' or ``it follows''?
\item If proving equality of sets, were both inclusions shown?
\item If proving a universal statement, was the object arbitrary?
\item If proving an existential statement, was a witness produced?
\item If proving by induction, were the base case, induction hypothesis, and induction step all explicit?
\end{enumerate}

\section{Feedback Form}

\begin{verbatim}
Claim:

My proof:

Where I feel uncertain:

Definitions I used:

Strategy I intended:

Feedback requested:
[ ] Check correctness
[ ] Find the first broken line
[ ] Suggest a proof strategy
[ ] Make the proof more elegant
[ ] Give a hint, not a full solution
\end{verbatim}

\end{document}
